\chapter{User Interface}
\label{ch:ui}

\section{Introduction}
The interface is where the system stops being a pipeline and becomes a public
service. This chapter describes the goals that governed its design, the visual
system it uses, and each of its four views. The figures in this chapter are
screenshots of the deployed application taken from the live public deployment, so
the values shown - the forecast, the leaderboard, the drift table - are the
system's genuine output rather than mock-ups.

\section{Interface Design Goals}
Five goals shaped the interface.

\textbf{Answer the question in the first screen.} A person opening the page wants
to know whether the air is safe. The current index, its health category and its
advice therefore appear immediately, in a card whose colour is itself the answer,
above the chart rather than below it.

\textbf{Never present a bare number.} In accordance with NFR-4, every forecast
carries a category name, an uncertainty band and, where one could be computed, an
explanation. A number alone invites either misplaced confidence or dismissal.

\textbf{Use the standard health semantics.} The colour scale, the category names
and the advisory wording are the EPA's~\cite{ref:epa}, not the project's. A user
who has seen an air-quality colour scale elsewhere should not have to learn a new
one, and the legend is present on the page so that the mapping is never implicit.

\textbf{Publish the system's own record.} Three of the four tabs are given over
to the model rather than the forecast: its promotion history, its accuracy by
horizon and its current drift status. This is unusual for a consumer-facing page
and it is deliberate, in accordance with principle P6.

\textbf{Degrade rather than break.} When the backend is unreachable the interface
says so specifically; when it is running without a backend at all, the two
capabilities that require a live model are removed from the navigation rather than
offered and failing.

\section{The Visual System}
The application uses a custom dark design system, developed for this project
under the name \emph{Aurora}. It is built on a near-black canvas with a slowly
drifting, heavily blurred aurora of cyan, teal, violet and sky gradients behind
the content, and translucent glass panels above it. The signature accent is a
cyan-to-teal gradient, chosen to sit well beside the black-and-white 10Pearls
mark, which appears as a circular badge in the header and as a wordmark in the
footer.

Three decisions in the visual system are functional rather than decorative. The
dark canvas maximises the contrast of the six EPA category colours, which are
saturated and are the most important information on the page; on a white canvas
the yellow and orange bands in particular lose legibility. The panels are
separated by elevation and translucency rather than by heavy borders, which keeps
the eye on the data. And the accent colour is deliberately not one of the six
category colours, so that interface chrome can never be mistaken for a health
signal.

An animated introduction holds the application's mark on screen for a short
minimum duration while the forecast document loads, then fades out once both the
minimum time has elapsed and the data has arrived. This turns the unavoidable
initial fetch into a deliberate moment rather than an empty page.

\section{The Forecast View}
Figure~\ref{fig:ui-forecast} shows the default view.

\fig{ui-forecast}{0.94}{The Forecast view of the deployed
application.}{fig:ui-forecast}

The view is arranged in two columns. The main column carries, from the top: the
hazard banner, shown only when the severity is not none; the current-conditions
card; the forecast chart; and the explanation card. The side column carries the
map, the category legend and the advisor panel.

The \textbf{hazard banner} is the highest-priority element on the page and is
present in the figure, reporting a peak of 200 in the \emph{unhealthy} band,
naming the hour at which the curve first crosses the threshold, and giving the
standard advice for that category. It appears above everything else because a
user who reads nothing else on the page should still receive the warning.

The \textbf{current card} states the observed index as a large numeral against the
colour of its category, names the category, repeats the health advice in full, and
adds the province and the peak expected over the next three days. Colour is used
as a redundant encoding here rather than as the sole one: the category is always
named in text, so the card remains readable to a colour-blind user.

The \textbf{forecast chart} draws the 72-hour curve with the 80\% prediction
interval as a shaded band behind it, and the category thresholds as faint
horizontal rules, so the reader can see at a glance which band the curve is in and
when it crosses. The band is labelled in the panel header rather than left to be
inferred. A toggle switches between the hourly view and a three-day aggregation
for a user who wants the summary rather than the detail.

The \textbf{map} places every city on a Pakistan basemap as a marker coloured by
its current category, with the selected city highlighted. It serves two purposes:
it allows a city to be selected by pointing rather than by finding it in a list,
and it makes the geographic pattern visible - in the figure, the concentration of
red markers across the Punjab plain against the yellow of the coast and the north
is the exploratory finding of Section~\ref{sec:design-eda} rendered as a picture.

The \textbf{explanation card} presents the sentence generated from the Shapley
attributions for the peak hour, followed by a signed horizontal bar for each of
the five leading drivers, and a note stating that the values are feature
contributions in index points for the peak forecast hour. In the figure the
current index contributes $+26$, the time of day $+13.4$, the season $-6.6$, the
six-hour trend $+6.4$ and the temperature $+5.0$ against a base value of 132 -
which tells a reader not merely that the air will be bad but that it will be bad
mainly because it is already bad and the diurnal cycle is working against it.

The \textbf{advisor panel} offers three suggested questions and a free-text field.
Suggestions are shown because a blank conversational field gives a user no
indication of what the system can answer.

\section{The Model Evaluation View}
Figure~\ref{fig:ui-eval} shows the evaluation view.

\fig{ui-eval}{0.94}{The Model Evaluation view.}{fig:ui-eval}

Three panels present the model's own record. The \textbf{leaderboard} names the
current champion with its version, family and metrics, and lists every training
run in reverse chronological order with its outcome. The refused run discussed in
Section~\ref{sec:impl-gate} is visible at the top of the table, marked
\emph{rejected}, immediately above the identical run marked \emph{champion}. That
a public interface shows a model that was refused is the point: it is evidence
that the gate is real.

The \textbf{error-by-horizon} chart plots the champion's RMSE against the
persistence baseline's at each of the ten modelled lead times. Two things are
legible at once: that error grows steeply with lead time, and that the champion is
below the baseline across almost the whole window. Publishing the first of these
is what makes the second credible.

The \textbf{walk-forward table} lists the five rolling folds with their training
size, test window start, RMSE and $R^2$, and states the mean. It is included
because a single validation number invites the suspicion that the split was
fortunate.

\section{The Monitoring View}
Figure~\ref{fig:ui-monitoring} shows the monitoring view.

\fig{ui-monitoring}{0.94}{The Monitoring view.}{fig:ui-monitoring}

The drift panel states the overall status and names the worst-drifting feature,
then tabulates the Population Stability Index for each monitored feature with a
coloured status badge. The three thresholds are stated in the panel's subtitle, so
that a reader can interpret a number without knowing the measure beforehand. In
the figure the overall status is significant, driven by temperature; particulate
matter at ten micrometres and wind speed remain stable.

Presenting an alarming word like \emph{significant} on a public page is a
deliberate choice, and the second panel is what makes it honest. Forecast error
tracking reports the realised accuracy of past forecasts once the hours they
described have arrived; in the figure it correctly reports that forecasts have
been logged and are waiting for observations rather than displaying a fabricated
zero. Read together, the two panels say what is true: the inputs have moved
seasonally, and whether that has cost any accuracy is a question the system will
answer with data rather than assumption.

\section{The What-If View}
Figure~\ref{fig:ui-whatif} shows the what-if view after a scenario has been run.

\fig{ui-whatif}{0.94}{The What-If view, showing a scenario result.}{fig:ui-whatif}

Five sliders expose the drivers described in Section~\ref{sec:impl-whatif}, each
initialised to the current real value for the selected city so that the user
begins from reality rather than from an arbitrary midpoint. Running the simulation
calls the live model twice and presents the baseline and the scenario side by
side, each as a large numeral on its category colour with its category named
beneath, and states the change between them.

The figure shows a heavy-pollution, still-air scenario for Lahore: particulate
matter raised to 280 micrograms per cubic metre, the current index raised to 310
and the wind reduced to 2 kilometres per hour. The model moves the 24-hour
forecast from 38, in the \emph{good} band, to 159, in the \emph{unhealthy} band -
a change of $+121$ index points. The panel's subtitle states plainly that this is
a model simulation and not causal proof, satisfying NFR-6.

The tab is present only when a live backend is configured, because the simulation
requires the model to be invoked; in the static serving mode it is removed from
the navigation entirely.

\section{Interaction and Responsiveness}
Navigation is a single row of tabs, and the selected city is shared across the
Forecast and What-If views, so a user who has chosen Multan does not have to
choose it again. The header states when the displayed forecast was generated,
which is essential for a product whose data is refreshed on a schedule: a stale
page should be identifiable as stale.

The layout is fluid. On a wide viewport the two columns sit side by side; on a
narrow one they stack, with the main column first, so that the current conditions
and the chart precede the map and the legend on a phone. Charts resize with their
container rather than being fixed, and long tables scroll horizontally within
their panel rather than forcing the page to scroll.

Three states are handled explicitly rather than being left to chance. While the
forecast document is loading, the application shows its mark, a spinner and a
caption. If the fetch fails, it shows a notice naming the failure and, for a
developer running locally, the command that starts the backend. If a selected city
is absent from the payload, the application falls back to the first city present
rather than rendering an empty view.

\section{Summary of User Interface}
This chapter described the interface through which the system reaches the public:
its five design goals, the dark Aurora visual system and the functional reasons
behind it, and its four views. The Forecast view answers the user's question
immediately and supports it with an uncertainty band, a map, a legend, an
explanation and an advisor. The Model Evaluation and Monitoring views publish the
system's own record - including a refused model and an unflattering drift status
- because a public-health signal that cannot be checked will not be trusted. The
What-If view lets a user interrogate the model directly. The next chapter reports
the testing and the empirical evaluation of everything described so far.
